\section{Grafos}

\subsection{BFS 0-1}

Usado em grafos cujas arestas possuem peso apenas $0$ ou $1$.

Mantém os vértices em uma \texttt{deque}:
\begin{itemize}
    \item peso $0$: adiciona no início;
    \item peso $1$: adiciona no final.
\end{itemize}

\textbf{Complexidade:} $O(V + E)$.

\begin{lstlisting}
struct Edge {
    int destination;
    int weight;
};

vector<long long> zeroOneBFS(const vector<vector<Edge>>& graph, int origin) {
    int vertexCount = graph.size();

    vector<long long> distance(vertexCount, INF);

    deque<int> vertices;

    distance[origin] = 0;
    vertices.push_front(origin);

    while (!vertices.empty()) {
        int currentVertex = vertices.front();
        vertices.pop_front();

        for (const Edge& edge : graph[currentVertex]) {
            long long newDistance = distance[currentVertex] + edge.weight;

            if (newDistance >= distance[edge.destination])
                continue;

            distance[edge.destination] = newDistance;

            // Arestas de peso 0 devem ser processadas antes das arestas de peso 1.
            if (edge.weight == 0)
                vertices.push_front(edge.destination);
            else
                vertices.push_back(edge.destination);
        }
    }

    return distance;
}
\end{lstlisting}


\subsection{BFS Multi-Estado}

Usado quando chegar ao mesmo vértice com estados diferentes
produz situações diferentes.

O princípio é tratar $(v, state)$ como um vértice diferente.

\textbf{Complexidade:} $O(VS + ES)$, onde $S$ é o número de estados.

\begin{lstlisting}
struct State {
    int vertex;
    int state;
};

vector<vector<int>> multiStateBFS(const vector<vector<int>>& graph, int origin, int stateCount) {
    int vertexCount = graph.size();

    vector<vector<int>> distance(vertexCount, vector<int>(stateCount, -1));

    queue<State> states;

    // Estado inicial.
    distance[origin][0] = 0;
    states.push({origin, 0});

    while (!states.empty()) {
        State current = states.front();
        states.pop();

        int currentDistance = distance[current.vertex][current.state];

        for (int neighbor : graph[current.vertex]) {
            // Exemplo: transicao que mantem o estado.
            int nextState = current.state;

            if (distance[neighbor][nextState] != -1)
                continue;

            distance[neighbor][nextState] = currentDistance + 1;

            states.push({neighbor, nextState});
        }
    }

    return distance;
}
\end{lstlisting}

O trecho que determina \texttt{nextState} deve ser adaptado ao problema.

A dimensão do vetor \texttt{distance} representa o espaço de estados.


\subsection{Dijkstra Multi-Estado}

Extensão do Dijkstra em que o vértice do grafo é combinado
com um estado adicional.

O estado completo é $(v, state)$.

\textbf{Complexidade:} $O((VS + ES)\log(VS))$.

\begin{lstlisting}
struct Edge {
    int destination;
    long long weight;
};

struct State {
    long long distance;
    int vertex;
    int state;

    bool operator>(const State& other) const {
        return distance > other.distance;
    }
};

vector<vector<long long>> multiDijkstra(const vector<vector<Edge>>& graph, int origin, int stateCount) {
    int vertexCount = graph.size();

    vector<vector<long long>> distance(vertexCount, vector<long long>(stateCount, INF));

    priority_queue<
        State,
        vector<State>,
        greater<State>
    > priorityQueue;

    distance[origin][0] = 0;

    priorityQueue.push({0, origin, 0});

    while (!priorityQueue.empty()) {
        State current = priorityQueue.top();
        priorityQueue.pop();

        // Ignora entradas antigas da priority_queue.
        if (current.distance != distance[current.vertex][current.state])
            continue;

        for (const Edge& edge : graph[current.vertex]) {
            // Adapte esta transicao conforme o problema.
            int nextState = current.state;

            long long newDistance = current.distance + edge.weight;

            if (newDistance >= distance[edge.destination][nextState])
                continue;

            distance[edge.destination][nextState] = newDistance;

            priorityQueue.push({newDistance, edge.destination, nextState});
        }
    }

    return distance;
}
\end{lstlisting}

O ponto mais importante é lembrar que:
\[
\boxed{\text{vértice real} \rightarrow (v,state)}
\]

Cada combinação é tratada como um vértice independente.


\subsection{MST --- Kruskal}

Encontra uma MST usando as arestas ordenadas por peso usando DSU para evitar ciclos.

\textbf{Complexidade:} $O(E\log E)$.

\begin{lstlisting}
using Edge = tuple<long long, int, int>;

long long kruskal(int vertexCount, vector<Edge> edges) {
    priority_queue<
        Edge,
        vector<Edge>,
        greater<Edge>
    > pq;

    for (const auto& [weight, firstVertex, secondVertex] : edges)
        pq.push({weight, firstVertex, secondVertex});

    DSU disjointSet(vertexCount);

    long long totalWeight = 0;
    int selectedEdges = 0;

    while (!pq.empty()) {
        auto [weight, firstVertex, secondVertex] = pq.top();
        pq.pop();

        // Ignora arestas que formariam um ciclo.
        if (!disjointSet.unite(firstVertex, secondVertex))
            continue;

        totalWeight += weight;
        selectedEdges++;

        if (selectedEdges == vertexCount - 1)
            break;
    }

    // Se nao houver arvore geradora, o grafo nao eh conexo.
    if (selectedEdges != vertexCount - 1)
        return -1;

    return totalWeight;
}
\end{lstlisting}


\subsection{MST --- Prim}

Adequado para grafos densos representados por matriz de adjacência.

A cada iteração, adiciona o vértice ainda não utilizado com a menor aresta conhecida para a MST.

\textbf{Complexidade:} $O(V^2)$.

\begin{lstlisting}
long long prim(const vector<vector<long long>>& graph) {
    int vertexCount = graph.size();

    vector<long long> minimumEdge(vertexCount, INF);

    vector<bool> used(vertexCount, false);

    long long totalWeight = 0;

    // Comeca a MST pelo vertice 0.
    minimumEdge[0] = 0;

    for (int iteration = 0; iteration < vertexCount; iteration++) {
        int currentVertex = -1;

        // Escolhe o vertice nao utilizado com a menor aresta para a MST.
        for (int vertex = 0; vertex < vertexCount; vertex++)
            if (!used[vertex] && (currentVertex == -1 || 
                minimumEdge[vertex] < minimumEdge[currentVertex]))
                currentVertex = vertex;

        // O grafo nao eh conexo.
        if (currentVertex == -1 || minimumEdge[currentVertex] == INF)
            return -1;

        used[currentVertex] = true;
        totalWeight += minimumEdge[currentVertex];

        // Atualiza a menor aresta que conecta cada vertice restante a MST.
        for (int neighbor = 0; neighbor < vertexCount; neighbor++)
            if (!used[neighbor] && graph[currentVertex][neighbor] < minimumEdge[neighbor])
                minimumEdge[neighbor] = graph[currentVertex][neighbor];
    }

    return totalWeight;
}
\end{lstlisting}


\subsection{Diâmetro de uma Árvore}

O diâmetro é o maior caminho entre dois vértices da árvore.

Em uma árvore não ponderada:

\begin{enumerate}
    \item Faça BFS/DFS a partir de qualquer vértice.
    \item Pegue o vértice mais distante $A$.
    \item Faça outra BFS/DFS a partir de $A$.
    \item A maior distância encontrada é o diâmetro.
\end{enumerate}

\textbf{Complexidade:} $O(V)$.

\begin{lstlisting}
pair<int, int> bfsFarthest(const vector<vector<int>>& tree, int origin) {
    int vertexCount = tree.size();

    vector<int> distance(vertexCount, -1);

    queue<int> vertices;

    distance[origin] = 0;
    vertices.push(origin);

    int farthestVertex = origin;

    while (!vertices.empty()) {
        int currentVertex = vertices.front();
        vertices.pop();

        if (distance[currentVertex] > distance[farthestVertex])
            farthestVertex = currentVertex;

        for (int neighbor : tree[currentVertex]) {
            if (distance[neighbor] != -1)
                continue;

            distance[neighbor] = distance[currentVertex] + 1;

            vertices.push(neighbor);
        }
    }

    return {farthestVertex, distance[farthestVertex]};
}

int treeDiameter(const vector<vector<int>>& tree) {
    // O vertice mais distante de qualquer vertice eh uma extremidade do diametro.
    auto [endpoint, unusedDistance] = bfsFarthest(tree, 0);

    // A maior distancia a partir dessa extremidade eh exatamente o diametro.
    auto [otherEndpoint, diameter] = bfsFarthest(tree, endpoint);

    return diameter;
}
\end{lstlisting}


\subsection{Grafo Bipartido}

Um grafo é bipartido se seus vértices podem ser divididos em dois conjuntos, 
de modo que nenhuma aresta conecte vértices do mesmo conjunto.

Pode ser verificado colorindo os vértices durante uma BFS/DFS.

\textbf{Complexidade:} $O(V + E)$.

\begin{lstlisting}
bool isBipartite(const vector<vector<int>>& graph) {
    int vertexCount = graph.size();

    // -1  = nao colorido
    // 0/1 = duas cores possiveis
    vector<int> color(vertexCount, -1);

    queue<int> vertices;

    for (int startVertex = 0; startVertex < vertexCount; startVertex++) {
        if (color[startVertex] != -1)
            continue;

        color[startVertex] = 0;
        vertices.push(startVertex);

        while (!vertices.empty()) {
            int currentVertex = vertices.front();
            vertices.pop();

            for (int neighbor : graph[currentVertex]) {
                if (color[neighbor] == -1) {
                    // Vertices adjacentes precisam receber cores diferentes.
                    color[neighbor] = color[currentVertex] ^ 1;

                    vertices.push(neighbor);
                }
                else if (color[neighbor] == color[currentVertex]) {
                    // Uma aresta entre vertices da mesma cor prova que nao eh bipartido.
                    return false;
                }
            }
        }
    }

    return true;
}
\end{lstlisting}


\subsection{Bellman-Ford}

Encontra menores caminhos a partir de uma origem em grafos com pesos negativos.

Também detecta ciclos negativos alcançáveis a partir da origem.

\textbf{Complexidade:} $O(VE)$.

\begin{lstlisting}
struct Edge {
    int origin;
    int destination;
    long long weight;
};

pair<vector<long long>, bool> bellmanFord(int vertexCount, const vector<Edge>& edges, int origin) {
    vector<long long> distance(vertexCount, INF);

    distance[origin] = 0;

    // O menor caminho simples possui no maximo V - 1 arestas.
    for (int iteration = 0; iteration < vertexCount - 1; iteration++) {
        bool changed = false;

        for (const Edge& edge : edges) {
            if (distance[edge.origin] == INF)
                continue;

            long long newDistance = distance[edge.origin] + edge.weight;

            if (newDistance < distance[edge.destination]) {
                distance[edge.destination] = newDistance;
                changed = true;
            }
        }

        // Nenhuma distancia mudou: ja convergiu.
        if (!changed)
            break;
    }

    // Se ainda for possivel relaxar uma aresta, existe um ciclo negativo alcancavel.
    for (const Edge& edge : edges) {
        if (distance[edge.origin] == INF)
            continue;

        if (distance[edge.origin] + edge.weight < distance[edge.destination])
            return {distance, true};
    }

    return {distance, false};
}
\end{lstlisting}


\subsection{Floyd-Warshall}

Encontra os menores caminhos entre \textbf{todos os pares}
de vértices.

Também funciona com arestas negativas, desde que não existam
ciclos negativos.

\textbf{Complexidade:} $O(V^3)$.

\textbf{Memória:} $O(V^2)$.

\begin{lstlisting}
void floydWarshall(vector<vector<long long>>& distance) {
    int vertexCount = distance.size();

    for (int intermediate = 0; intermediate < vertexCount; intermediate++) {
        for (int origin = 0; origin < vertexCount; origin++) {
            if (distance[origin][intermediate] == INF)
                continue;

            for (int destination = 0; destination < vertexCount; destination++) {
                if (distance[intermediate][destination] == INF)
                    continue;

                distance[origin][destination] = 
                    min(distance[origin][destination], 
                        distance[origin][intermediate] + distance[intermediate][destination]);
            }
        }
    }
}
\end{lstlisting}

Inicialização:

\begin{lstlisting}
vector<vector<long long>> distance(vertexCount, vector<long long>(vertexCount, INF));

for (int vertex = 0; vertex < vertexCount; vertex++) 
    distance[vertex][vertex] = 0;

for (const Edge& edge : edges) {
    distance[edge.origin][edge.destination] =
        min(distance[edge.origin][edge.destination], edge.weight);
}
\end{lstlisting}

Após o algoritmo, se:
\[
distance[v][v] < 0
\]

para algum vértice $v$, existe um ciclo negativo.